\subsection{The inseparable Albanese surface branch}
\label{subsec:smallest-endpoint-albanese-geometry}

We now impose the Albanese hypotheses.  Throughout this subsection,
$k$ is algebraically closed of characteristic two, $X$ is a smooth
integral projective surface, and $C\subset X$ is a smooth integral
ample divisor.  Let $M\ne\OO_X$ be algebraically trivial of exact
order two, with $M|_C\simeq\OO_C$.  Assume that its logarithmic form
$\omega$ has exact order two along $C$ and that
\[
 d=C^2,\qquad 2g(C)-2=3d.
\]
Choose a base point on $C$, let
$\alpha:X\to A=\operatorname{Alb}(X)$ be the pointed Albanese map,
and suppose that $\dim\alpha(X)=2$.  Put $q_A=\dim A\ge2$.
Since an algebraically trivial line bundle has a square root,
Theorem~\ref{thm:smallest-endpoint-global-structure} applies.
We retain its notation
\[
 N=\OO_C(C),\qquad
 \mathcal T=\omega_C\otimes N^{-3},\qquad
 \pi:V\longrightarrow X,\qquad
 \theta=\omega/s_C^2,\qquad
 \mathcal F=\omega_X^{-1}(2C).
\]
In particular, $\mathcal T$ has exact order two.

\begin{proposition}[The canonical quotient and its Albanese variety]
\label{prop:smallest-endpoint-albanese-quotient}
The Albanese differential has generic rank exactly one, and
\[
 T_{X/A}=\mathcal F.
\]
Let $q:X\to Y$ be the quotient by this foliation and write
$\alpha=hq$.  Then $Y$ is normal, $q$ is finite and purely inseparable
of degree two, and
\begin{equation}\label{eq:smallest-endpoint-albanese-quotient}
 Y\simeq V^{(2)},\qquad q\pi=F_{V/k}.
\end{equation}
Here $F_{V/k}:V\to V^{(2)}$ denotes relative Frobenius.
The map $h:Y\to A$ is the Albanese map of $Y$.
Consequently
\begin{equation}\label{eq:smallest-endpoint-albanese-frobenius}
 \operatorname{Alb}(V)^{(2)}\simeq A,
\end{equation}
and, under this identification, the homomorphism induced by
$\pi$ is the relative Frobenius of $\operatorname{Alb}(V)$.
There is no assertion that $h$ is generically separable.
\end{proposition}

\begin{proof}
Picard--Albanese duality identifies $\alpha^*$ on $\Pic^0(k)$
with the algebraically trivial line bundles on $X$
\cite[Expos\'e~236, Corollary~3.2 and Theorem~3.3(iii),
pp.~236--238]{FGAPicard}.
Thus $M=\alpha^*P$ for a unique nontrivial order-two line bundle
$P$ on $A$.  Its logarithmic form $\omega_P$ is a nonzero invariant
one-form and satisfies $\alpha^*\omega_P=\omega$.
Nonzeroness follows from logarithmic injectivity: a zero derivative
would give local unit square roots and trivialize $P$.
The global collinearity in
Theorem~\ref{thm:smallest-endpoint-global-structure} gives the upper
bound one for the rank of $d\alpha$, while $\omega$ gives the lower
bound.  Invariant forms generate $\Omega_A^1$, so their common
annihilator is exactly $\mathcal F=T_{X/A}$.

This is the canonical degree-two foliation associated with an
inseparable Albanese map with surface image
\cite[Lemma~2.22 and Remark~2.23, p.~13]{ContiPositiveSeveriComparison}.
We verify the normal quotient and its identification explicitly.
Put $K=k(X)$, and write a rational torsor equation as $z^2=u$.
Since $u\notin K^2$, it can be completed to a $2$-basis $u,v$ of $K$.
The constants of a nonzero derivation annihilating $u$ are
$K^2(u)$: using the basis $1,u,v,uv$ over $K^2$, this is the
kernel of differentiation with respect to $v$.
On the other hand,
\[
 k(V)^2=K^2(u).
\]

Let $R$ be a normal affine coordinate ring in $X$, and let
$\widetilde R$ be its integral closure in $k(V)$.
Every $b\in\widetilde R$ satisfies $b^2\in R$, because $b^2$
belongs to $K$ and is integral over $R$.
Frobenius on $V$ therefore factors through a map
$q':X\to V^{(2)}$.
This map is finite: locally $R$ is finite over $R^2$, hence over
the intermediate Frobenius image ring $\widetilde R^2$.
It has degree two by the field calculation above.
Moreover,
\[
 R\cap k(V)^2=\widetilde R^2.
\]
Indeed, if $r=b^2\in R$ with $b\in k(V)$, then $b$ is integral over
$R$ and belongs to $\widetilde R$.
Thus the structure sheaf of $V^{(2)}$ is precisely the invariant
subsheaf for the foliation.  These identifications use the usual
Frobenius-twist convention over the perfect field $k$.
They identify $q'$ with the normal quotient $q$.
Equivalently, both normal quotients are the integral closures of
the Frobenius base ring in the same intermediate field.
This proves \eqref{eq:smallest-endpoint-albanese-quotient}.
Since the foliation annihilates the pullbacks of local functions
on $A$, the invariant-sheaf description also gives the factor $h$.

For any pointed map $\psi:Y\to A_0$ to an abelian variety,
universality of $\alpha$ gives a unique homomorphism $v:A\to A_0$
such that
\[
 \psi q=v\alpha=vhq.
\]
Equality descends through $q$: it is dominant and $Y$ is reduced,
so its pullback on functions is injective, and the target is
separated.  Hence $\psi=vh$, proving the Albanese universal
property for $h$.
Formation of the Albanese commutes with base change in the cited
FGA theorem.  Applied to $Y\simeq V^{(2)}$, this gives
\eqref{eq:smallest-endpoint-albanese-frobenius}.
The map induced by $q$ is the identity under
$\operatorname{Alb}(Y)=A$, and $q\pi=F_{V/k}$.
The universal property of relative Frobenius now identifies the
map induced by $\pi$ with $F_{\operatorname{Alb}(V)/k}$.
This calculation does not remove any further inseparability
from $h$.
\end{proof}

\begin{proposition}[Invariant forms on the distinguished curve]
\label{prop:smallest-endpoint-albanese-forms}
Let
\[
 \mathcal W=\alpha^*H^0(A,\Omega_A^1)\subset H^0(X,\Omega_X^1).
\]
Then $\dim\mathcal W=q_A$, and tangential restriction to $C$ gives
\begin{equation}\label{eq:smallest-endpoint-albanese-restriction}
 \ker\bigl(\mathcal W\longrightarrow H^0(C,\omega_C)\bigr)
   =k\omega,\qquad
 \dim\operatorname{im}\bigl(\mathcal W\longrightarrow
 H^0(C,\omega_C)\bigr)=q_A-1.
\end{equation}
In particular, $C$ maps finitely and separably to the normalization
of $\alpha(C)$.
\end{proposition}

\begin{proof}
Pullback of invariant one-forms by the Albanese map is injective
for a smooth projective variety in positive characteristic
\cite[Lemma~1.3 and its proof, p.~195]{MehtaSrinivasAlbanese}.
Thus $\dim\mathcal W=q_A$.
All these forms are closed, and $\omega\in\mathcal W$.

Let $0\ne\beta\in\mathcal W$ have zero tangential restriction.
Its full restriction to $\Omega_X^1|_C$ lies in the conormal line
$N^{-1}$, which has negative degree, and is therefore zero.
Let $j\ge1$ be the exact vanishing order of $\beta$ along $C$.
After division by $C^j$, its leading coefficient has nonzero
tangential image: otherwise it would be a nonzero section of
$N^{-(j+1)}$, again of negative degree.
In local parameters $t,x$, write
\[
 \beta=t^j(a\,dt+b\,dx)
\]
with $b$ nonzero at the generic point of $C$.
The coefficient of $t^{j-1}dt\wedge dx$ in $d\beta$ is $jb$.
Closedness forces $j$ to be even, hence $j\ge2$.

By \eqref{eq:smallest-endpoint-global-forms},
$\beta=\theta s$ for a section $s$ of $\OO_X(2C)$.
The divided form $\theta$ has exact order zero along $C$,
so $s$ is divisible by $s_C^2$.
The quotient is a global regular function on $X$, hence a constant.
Thus $\beta\in k\omega$; the reverse inclusion is immediate.
This proves \eqref{eq:smallest-endpoint-albanese-restriction}.
Its image has dimension $q_A-1>0$, so $\alpha|_C$ has nonzero
differential.  The resulting nonconstant map to the normalized
image curve is finite and separable.
This assertion concerns the curve map, not the surface map.
\end{proof}

\begin{proposition}[The curve divisor and Albanese dimension]
\label{prop:smallest-endpoint-albanese-dimension}
Let $R$ be the zero divisor of the tangential component of
$\theta|_C$, a section of $\omega_C\otimes N^{-2}$.
There is an effective divisor $E_C$ such that
\begin{equation}\label{eq:smallest-endpoint-albanese-curve-lines}
 R=2E_C,\qquad \deg E_C=d/2,\qquad
 N\simeq\OO_C(2E_C)\otimes\mathcal T,\qquad
 \omega_C\simeq\OO_C(6E_C).
\end{equation}
Let $R_\alpha$ be the common zero divisor of the nonzero image
space in \eqref{eq:smallest-endpoint-albanese-restriction}.
Then
\begin{equation}\label{eq:smallest-endpoint-albanese-dimension}
 d\le\deg R_\alpha\le3d-2q_A+4,\qquad
 2\le q_A\le d+2.
\end{equation}
\end{proposition}

\begin{proof}
The tangential component is nonzero.  Otherwise the nonzero leading
restriction of $\theta$ would lie in $N^{-3}$, which has negative
degree.
The local model
\eqref{eq:smallest-endpoint-global-local-model} gives the
tangential component as
\[
 \frac{d(A_i|_C)}{(b_i|_C)^2}.
\]
In a completed local parameter on $C$, the derivative of a power
series in characteristic two has only even powers.  Its zero order
is even, and the denominator is a square unit.
Consequently $R=2E_C$.
Its degree is $2g(C)-2-2d=d$, and
\[
 \OO_C(R)=\omega_C\otimes N^{-2}=N\otimes\mathcal T.
\]
Since $\mathcal T^2\simeq\OO_C$, these identities give
\eqref{eq:smallest-endpoint-albanese-curve-lines}.

Equation~\eqref{eq:smallest-endpoint-global-forms}, restricted
tangentially to $C$, shows that $R_\alpha\ge R$.
Put $H=\omega_C(-R_\alpha)$, and write
\[
 r=h^0(C,H),\qquad s=h^0(C,\omega_C\otimes H^{-1}).
\]
We have $r\ge q_A-1$ and
$s=h^0(C,\OO_C(R_\alpha))\ge1$.
We give the Clifford estimate needed here directly.
At a smooth point, choose bases of these two spaces adapted to
vanishing order, with strictly increasing orders
$a_1<\cdots<a_r$ and $b_1<\cdots<b_s$.
The products with orders
\[
 a_1+b_1<\cdots<a_r+b_1<a_r+b_2<\cdots<a_r+b_s
\]
are $r+s-1$ independent canonical sections.
Hence $r+s-1\le g(C)$.
Combining this with Riemann--Roch,
$r-s=\deg H+1-g(C)$, gives
$2r\le\deg H+2$.
Therefore
\[
 2q_A-2\le3d-\deg R_\alpha+2.
\]
Together with $\deg R_\alpha\ge d$, this proves
\eqref{eq:smallest-endpoint-albanese-dimension}.
\end{proof}

\begin{remark}
\label{rem:smallest-endpoint-albanese-zero-divisor-scope}
The divisor $R_\alpha$ records zeros of the differential of the map
to $A$.  It need not be the ramification divisor of the separable
map from $C$ to its normalized image: the map from that normalized
curve into $A$ can itself have zeros in its differential.
No identification of these two divisors is used.
\end{remark}

\begin{proposition}[The normalized image curve]
\label{prop:smallest-endpoint-albanese-image}
Let $B_0$ be the normalization of $\alpha(C)$, write
$f:C\to B_0$ for the finite separable map, and
$j:B_0\to A$ for the normalization map followed by inclusion.
For the line bundle $P$ in the proof of
Proposition~\ref{prop:smallest-endpoint-albanese-quotient},
\begin{equation}\label{eq:smallest-endpoint-albanese-image-torsion}
 j^*P\simeq\OO_{B_0}.
\end{equation}
The curve $j(B_0)$ generates $A$.  In particular, with $b=g(B_0)$,
one has $b\ge q_A\ge2$, and
\[
 \operatorname{Jac}(B_0)\longrightarrow A
\]
is surjective and inseparable.  Its pullback on invariant
cotangent spaces has one-dimensional kernel.
\end{proposition}

\begin{proof}
If $j^*P$ were nontrivial, its $\mu_2$-torsor over the smooth curve
$B_0$ would give a nontrivial purely inseparable quadratic extension
of $k(B_0)$.
Its pullback line to $C$ is trivial because $M|_C$ is trivial.
The pulled-back square trivialization differs from a square frame
by a constant; every constant in $k^*$ has a square root.
Thus the pulled-back torsor is trivial, placing its defining square
root in $k(C)$.  This contradicts separability of
$k(C)/k(B_0)$ and proves
\eqref{eq:smallest-endpoint-albanese-image-torsion}.

Let $A_0$ be the abelian subvariety generated by $j(B_0)$.
Our base point on $C$ makes this curve pass through the origin.
The composite $X\to A/A_0$ contracts $C$.
If it were nonconstant, choose a hyperplane section of a projective
embedding of its image which avoids the image of $C$ and meets
that positive-dimensional image.  Its pullback is a nonzero
effective divisor disjoint from $C$, contrary to ampleness.
Hence $\alpha(X)\subset A_0$, and the universal property of
$\alpha$ forces $A_0=A$.
The Jacobian homomorphism is therefore surjective, giving
$b\ge q_A\ge2$.

Pullback of regular one-forms by the separable curve map $f$ is
injective.  Thus the kernel of $j^*$ on invariant one-forms is
the kernel after pulling back further to $C$.
By \eqref{eq:smallest-endpoint-albanese-restriction} and injectivity
of $\alpha^*$, it is precisely $k\omega_P$.
Under the usual identification of invariant differentials on
$\operatorname{Jac}(B_0)$ with $H^0(B_0,\omega_{B_0})$,
this is the cotangent kernel of the Jacobian homomorphism.
Its differential has rank $q_A-1$, so the surjective homomorphism
is inseparable.
\end{proof}

\begin{remark}
\label{rem:smallest-endpoint-albanese-normalized-image-scope}
The normalized image curve $B_0$ is not thereby the base of a
morphism from $X$.  In particular, the curve-pullback excess
inequality cannot be applied to $B_0$ without an additional
construction of such a morphism and a pullback description of $M$.
\end{remark}

\begin{proposition}[The image divisor on an abelian surface]
\label{prop:smallest-endpoint-albanese-surface-defect}
Suppose $q_A=2$ and put $D_A=\alpha(C)\subset A$ with its reduced
integral structure.  Let $\rho:A'\to A$ be the $\mu_2$-isogeny
associated with $P$.
The map $j:B_0\to A$ lifts to $j':B_0\to A'$.
If $D'$ is its integral image and
$\epsilon=p_a(D')-b$, then $\epsilon\ge0$ and
\begin{equation}\label{eq:smallest-endpoint-albanese-surface-defect}
 \rho^*D_A=2D',\qquad
 p_a(D_A)-b=b-1+2\epsilon,\qquad
 D_A^2=4(b-1+\epsilon).
\end{equation}
Thus $D_A$ is singular, its normalization defect is at least $b-1$,
and $D_A^2$ is a positive multiple of four.
Equality in the defect bound holds exactly when $D'$ is smooth.
These are assertions about the image divisor, not a formula for
$C^2$.
\end{proposition}

\begin{proof}
The $\mu_2$-cover of the nontrivial line $P$ is integral.
It is smooth because its local defining equation has differential
the nonzero invariant form $\omega_P$, which has no zeros.
With a point over the origin it is an abelian variety and
$\rho$ is an isogeny; this is the abelian-cover construction in
the proof of \cite[Lemma~1.3, p.~195]{MehtaSrinivasAlbanese}.
Equation~\eqref{eq:smallest-endpoint-albanese-image-torsion},
with the same square-scalar adjustment as above, gives the lift
$j'$.

Since $j$ is birational onto $D_A$, the intermediate fields satisfy
\[
 k(D_A)\subset k(D')\subset k(B_0)=k(D_A).
\]
Thus $D'\to D_A$ is birational, and $B_0$ normalizes both curves.
The radicial degree-two isogeny has a unique divisor above $D_A$.
Its residue-field degree there is one, so the ramification index
is two and $\rho^*D_A=2D'$.
Finite intersection projection gives
\[
 2D_A^2=(\rho^*D_A)^2=4(D')^2.
\]
Adjunction on the two smooth abelian surfaces gives
\[
 p_a(D_A)=1+\frac{D_A^2}{2},\qquad
 p_a(D')=1+\frac{(D')^2}{2}.
\]
Substituting $p_a(D')=b+\epsilon$ yields
\eqref{eq:smallest-endpoint-albanese-surface-defect}.
The normalization defect $\epsilon$ is nonnegative.
It vanishes exactly when the integral curve $D'$ is normal,
equivalently smooth over the perfect field $k$.
Since $b\ge2$, the asserted strict singularity of $D_A$ follows.
\end{proof}

\begin{remark}[The image-curve pattern]
\label{rem:smallest-endpoint-albanese-image-pattern}
The image-divisor conditions alone are compatible with examples.
Take ordinary elliptic curves $E_1,E_2$ and the degree-two
$\mu_2$-isogeny
\[
 \rho=F_{E_1/k}\times1:
 A'=E_1\times E_2\longrightarrow E_1^{(2)}\times E_2=A.
\]
Choose a smooth sufficiently ample integral curve $D'\subset A'$
whose projection to $E_2$ has nonzero differential.
Such a choice follows from Bertini in a sufficiently ample linear
system separating first-order jets: impose at one smooth point
a tangent direction not killed by the second projection, and
choose a general member smooth elsewhere.
The restriction of $\rho$ to $D'$ is radicial and has nonzero
differential.  It is therefore birational onto its image $D_A$.
Its normalization kills the line $P$ defining $\rho$, and the
defect equality in
\eqref{eq:smallest-endpoint-albanese-surface-defect} holds with
$\epsilon=0$.

This constructs only the image-curve behavior.
Taking $X=A'$ would make $\rho^*P$ trivial, and $\rho$ would not be
the Albanese map of $A'$.
It is not a construction of a surface satisfying the endpoint
hypotheses.  The existence or exclusion of the full inseparable
Albanese endpoint configuration remains open.
\end{remark}
